% ============================================================
%  Section 5: Pre-training Strategy
% ============================================================

\section{Pre-training Strategy}
\label{sec:pretraining}

\bert{} is pre-trained via \emph{self-supervised learning} on large unlabelled
corpora.  Two complementary objectives are optimised jointly, allowing the model to
acquire both token-level and sentence-level linguistic understanding.

\subsection{Masked Language Modelling (MLM)}
\label{subsec:mlm}

The MLM objective is the principal innovation that distinguishes \bert{} from
left-to-right pre-training approaches.

\paragraph{Procedure.}  At each training step, $15\%$ of the WordPiece tokens in
an input sequence are selected uniformly at random.  Of these selected tokens:
\begin{itemize}
  \item $80\%$ are replaced by the special \mask{} token;
  \item $10\%$ are replaced by a randomly sampled vocabulary token; and
  \item $10\%$ are left unchanged.
\end{itemize}
The model is then trained to predict the \emph{original} token at each masked
position using its final hidden state, producing the probability:
\begin{equation}
  P\!\left(w_i \mid \mathbf{x}_{\setminus i}\right)
  = \softmax\!\left(\mathbf{W}_o\,\mathbf{h}_i^{(L)} + \mathbf{b}_o\right),
  \label{eq:mlm}
\end{equation}
where $\mathbf{h}_i^{(L)}$ is the top-layer hidden state for position $i$,
$\mathbf{W}_o \in \R^{|\mathcal{V}| \times d}$ is a learned output projection, and
$|\mathcal{V}|$ is the vocabulary size.

\paragraph{Significance.}  Because the hidden state for a masked token is computed
from its \emph{full bilateral context}, the model is forced to integrate information
from both directions simultaneously.  This is fundamentally different from
autoregressive LMs, where causally masked attention prevents right-to-left
information flow.

\paragraph{Training signal quality.}  The 10\% random replacement and 10\%
unchanged strategies are crucial: they prevent the model from learning that
\mask{} always indicates a position to predict, while still providing the majority
of gradient signal through masked positions.

\begin{table}[H]
  \centering
  \caption{Masked Language Modelling (MLM) token treatment strategy. For each of
           the 15\% selected tokens, the replacement is chosen probabilistically.}
  \label{tab:mlm_strategy}
  \renewcommand{\arraystretch}{1.4}
  \begin{tabular}{lcc}
    \toprule
    \textbf{Action} & \textbf{Proportion} & \textbf{Purpose} \\
    \midrule
    Replace with [MASK] & $80\%$ & Primary training signal \\
    Replace with random token & $10\%$ & Encoder robustness \\
    Keep unchanged & $10\%$ & Prevent [MASK] recognition \\
    \bottomrule
  \end{tabular}
\end{table}

\subsection{Next Sentence Prediction (NSP)}
\label{subsec:nsp}

Many downstream tasks --- notably natural language inference, question answering,
and reading comprehension --- require the model to reason about relationships
\emph{between} sentences.  To pre-train this capability, \bert{} introduces an
auxiliary binary classification task.

\paragraph{Construction.}  For each training example, a sentence pair~$(A, B)$ is
constructed as follows:
\begin{itemize}
  \item With probability $50\%$: $B$ is the sentence that actually follows $A$ in
        the source corpus (\texttt{IsNext}).
  \item With probability $50\%$: $B$ is a sentence sampled uniformly at random
        from the corpus (\texttt{NotNext}).
\end{itemize}
The final hidden state of the \cls{} token is fed into a binary classifier to
predict whether the pair is contiguous.

\paragraph{Limitations.}  Subsequent work~\cite{liu2019roberta} demonstrated that
the NSP objective provides limited benefit and can even hurt performance when the
training data is already richly informative.  RoBERTa, a later variant, omits NSP
entirely without degradation.

\subsection{Pre-training Corpora}
\label{subsec:data}

\bert{} is pre-trained on two large English corpora:
\begin{itemize}
  \item \textbf{BooksCorpus}~\cite{zhu2015aligning}: approximately 800M words of
        unpublished books, providing long contiguous passages essential for learning
        inter-sentence relations.
  \item \textbf{English Wikipedia}: approximately 2,500M words (text only;
        lists, tables, and headers excluded), providing factual and encyclopaedic
        knowledge.
\end{itemize}
In total, pre-training uses roughly 3.3B tokens and requires training for
$1\,000\,000$ steps with a batch size of 256 sequences, amounting to approximately
$4$ days on $64$ Cloud TPU v3 chips for \bert{}-Base.

\subsection{Key Observations}
\label{subsec:pretrain_obs}

\begin{enumerate}
  \item \textbf{Implicit linguistic supervision.}  No hand-labelled data is used;
        supervision is derived entirely from the structure of natural language text.
  \item \textbf{Bidirectionality.}  The MLM objective produces representations
        that encode context from both directions, unlike GPT-style left-to-right
        pre-training.
  \item \textbf{High transferability.}  The resulting representations generalise
        across tasks with orthogonal linguistic requirements, requiring only a
        small task-specific fine-tuning head.
\end{enumerate}
